\cleardoublepage%
\phantomsection%
\addcontentsline{toc}{chapter}{Abstract}
\chapter*{Abstract}
\markboth{Abstract}{Abstract}


Electronic support measures systems intercept and analyse radar signals to identify the emitters present and how they operate among other things.

After the receiver, each recording is traditionally reduced to a time-ordered stream of pulse descriptor words, in which pulses from several emitters are interleaved and each emitter may change its operating mode.


Operators need both groupings from that mixed list: which pulses belong to the same physical emitter, and which pulses share an operating mode.
Classical pipelines deinterleave first and then classify modes on the recovered trains.
In such a cascade the mode stage takes the deinterleaver's hard pulse assignments as given, with no notion of how reliable they are, so a split or contaminated train silently corrupts the input to mode recognition.
The two tasks also rely on overlapping pulse-level cues such as carrier frequency, pulse width, and pulse repetition pattern, which suggests that a representation shared between them, as in multi-task learning, may serve both better than two separately designed stages.
This thesis therefore treats the two partitions as a joint clustering problem.
A multi-task transformer encoder architecture is proposed, composed of a shared trunk and two task-specific branches maps a window of pulses to an emitter embedding and a mode embedding.
The branches are trained with two supervised contrastive losses that differ only in how positives are chosen: emitter identifiers versus mode labels.
At inference, density-based clustering (DBSCAN) is applied independently to the mode and emitter embeddings, without assigning names from a fixed catalogue.

We train and evaluate the encoder on a synthetic, fully labelled scenario corpus.
Each pulse carries a recording-local emitter identity and a mode from a catalogue of $19$ types shared across recordings.
Recordings are split into windows of $2000$ pulses, and scores are reported on $1169$ non-overlapping test windows.
We compare the joint objective with single-task training, and standard self-attention with rotary encoding of time of arrival and with a pairwise physical attention bias.

Standard self-attention already recovers both partitions on most windows, with mean emitter adjusted Rand index (ARI) $0.948$ and mean mode ARI $0.985$.
Encoding time of arrival in attention, trogh RoPE or a introduced pairwise physical attention bias, each raises emitter ARI to $1.000$ and removes the remaining failure tail.
Combining the two does not.
Under standard attention, without any encoding of time of arrival, a deinterleaving-only encoder outperforms the joint model on emitters (ARI $0.999$ versus $0.948$).
Once time is encoded in attention, the joint encoder narrowly outperforms the single-task vanilla models.
The joint objective is therefore a trade-off under standard attention, not a free extra head.
Held-out modulations and operational recordings are left to future work.

\vspace{1em}
\noindent\textbf{Keywords:} transformer, deep clustering, electronic intelligence, radar emitter identification, mode recognition, supervised contrastive learning.
